% =============================================================================
% example-technical-implementation.tex
% Style: latex-docs-technical-implementation
% Build: pdflatex via latexmk; minted-aware with listings fallback.
% =============================================================================
\documentclass[11pt]{article}
\usepackage{latex-docs-technical-implementation}

\setDocTitle{Tile Binner -- Implementation Guidance}
\setDocSubtitle{CPU-side per-frame light culling for the forward
                tile-binned pipeline}
\setDocVersion{v0.2}
\setDocStatus{Draft}
\setDocOwner{Engine Architecture}
\setDocAuthor{Jordan Suber}
\setDocDate{May 8, 2026}
\setDocSummary{%
  Reference implementation guidance demonstrating the implementation
  subtype: API contracts with structured Inputs / Outputs / Errors /
  Side effects, inline source-code references, and a numbered
  walkthrough.%
}

\begin{document}
\makedoctocpage

\section{Component overview}
The \Module{TileBinner} runs once per frame on the CPU. It consumes
the view-frustum and the active light set, and produces a per-tile
light list. The list is uploaded to GPU as a single buffer texture
and consumed by the forward shader stage. The component is
implemented in \CodeRef{src/render/tile\_binner.c}{1}{file header},
and is invoked by the frame graph at the start of the lighting pass.

\section{Public API}

\begin{apicontract}{tilebinner\_run}{C function}
  \Inputs A populated \TypeName{tb\_input\_t}: view-projection matrix,
          tile dimensions, viewport size, and the active
          \TypeName{light\_list\_t}.
  \Outputs A \TypeName{tb\_output\_t}: a flat \texttt{uint32\_t} buffer
          of per-tile light index lists, prefixed by per-tile counts.
  \Errors  Returns \TypeName{TB\_ERR\_OOM} if the per-tile arena cannot
           grow; \TypeName{TB\_ERR\_INVAL} if the input is malformed.
           Both errors are non-fatal; the caller may render a previous
           frame's bins.
  \SideEffects Allocates from the per-frame arena passed in
           \TypeName{tb\_input\_t.arena}; never allocates from the
           global heap. Logs a debug counter via
           \Function{rdr\_metrics\_inc} when binning exceeds budget.
\end{apicontract}

\begin{apicontract}{tilebinner\_upload}{C function}
  \Inputs A \TypeName{tb\_output\_t} produced by
          \Function{tilebinner\_run} and a \TypeName{tb\_gpu\_handles\_t}.
  \Outputs Updates the buffer-texture handle in
          \TypeName{tb\_gpu\_handles\_t.light\_indices}; safe to call
          from the render thread.
  \Errors  Returns \TypeName{TB\_ERR\_GPU} if the buffer-texture
           upload fails; the caller should retry next frame.
  \SideEffects None on the CPU side beyond the GPU upload.
\end{apicontract}

\section{Walkthrough}

\begin{walkthrough}{One frame of binning}
  \WStep{The frame graph calls \Function{tilebinner\_run} immediately
         after the visibility query completes; see
         \CodeRef{src/render/frame\_graph.c}{142}{the binning hook}.}
  \WStep{For each light, the binner computes the screen-space bounding
         rectangle of the light's volume of effect; see
         \CodeRef{src/render/tile\_binner.c}{38}{\Function{light\_screen\_bbox}}.}
  \WStep{For each tile the rectangle overlaps, the binner appends the
         light's index to that tile's per-frame list. The list lives in
         the per-frame arena; no heap allocation in the hot path.}
  \WStep{When binning is complete, \Function{tilebinner\_compact}
         flattens the per-tile lists into the contiguous output buffer
         by writing the per-tile counts as a prefix; see
         \CodeRef{src/render/tile\_binner.c}{211}{the compaction step}.}
  \WStep{The frame graph calls \Function{tilebinner\_upload}, which
         binds the buffer texture and uploads it as a single
         \Function{glBufferSubData} call. \Constraint{This is the only
         CPU-to-GPU upload path the binner uses}; do not introduce
         additional uploads without revisiting the per-frame budget.}
\end{walkthrough}

\section{Performance notes}

\begin{infobox}[Where the time goes]
On the reference scene with 128 active lights, the binner takes
0.34~ms (\Function{tilebinner\_run}) plus 0.06~ms for the GPU upload
(\Function{tilebinner\_upload}). Per-tile light iteration in the
fragment shader costs an additional 0.5~ms on the reference GPU.
Total tile-binning overhead is therefore around 1~ms of a 16~ms frame
budget.
\end{infobox}

\section{What not to do}

\begin{cautionbox}[Common pitfalls]
\begin{itemize}
  \item Do \emph{not} allocate per-tile lists from the global heap; the
        per-frame arena is mandatory.
  \item Do \emph{not} run the binner on the render thread; the CPU
        cost belongs to the simulation thread per the frame-graph
        contract.
  \item Do \emph{not} cache binner output across frames. Determinism is
        a requirement of the design (REQ-RP-003); a cached bin breaks
        the visual regression suite.
\end{itemize}
\end{cautionbox}

\end{document}



